\documentclass[11pt,a4paper]{article}

\usepackage[utf8]{inputenc}
\usepackage[T1]{fontenc}
\usepackage{lmodern}
\usepackage{amsmath,amssymb,amsfonts,bm}
\usepackage{graphicx}
\usepackage{booktabs}
\usepackage{hyperref}
\usepackage{url}
\usepackage[margin=1in]{geometry}
\usepackage{enumitem}
\setlength{\emergencystretch}{2em}
\usepackage{multirow}

\title{Sizing and Capacity Boundaries of\\Spatial Convolutional Decoders\\for Surface-Code Quantum Error Correction}

\author{Brian Guarraci\\
  Galatheus\\
  \texttt{brian@galatheus.com}}

\date{October 2026 revision}

\begin{document}
\maketitle

% ====================================================================
\begin{abstract}
We study direct neural decoding for the rotated surface code by embedding
detector events as a three-dimensional volume: two spatial dimensions from
the code lattice and one temporal dimension from syndrome measurement
rounds. A residual 3D convolutional network operates on this volume and is
conditioned on the physical error rate, so one model covers a range of
sub-threshold operating points. The purpose is not to claim
state-of-the-art MWPM beating---hybrid and graph-based neural decoders
already establish that---but to characterise, with full multi-seed
coverage, \emph{how a pure spatial decoder performs under specified
architecture and training budgets}.
Across three seeds and a disjoint test split, the decoder beats
minimum-weight perfect matching (MWPM) at distances $d=3,5,7$ (relative
improvement $+11$ to $+26\%$ at physical error rate $p=0.003$), reaches a
noisy boundary at $d=9$ at the selected knee (tied with MWPM, large
seed-to-seed variance)---which a $5\times$-data follow-up resolves, beating MWPM
by $+36$ to $+57\%$ (three seeds) and showing the $d=9$ boundary is
under-resourcing rather than a fundamental limit---and underperforms at
$d=11$ (well below MWPM at the selected knee). We size the network per distance by an empirical
$\textrm{width}\times\textrm{depth}$ knee, and compare the classic ResNet
block against a Leela-Chess-Zero-style block (``LC0''): block design is
\emph{not} second-order but distance-dependent---LC0 is neutral at
$d=3,5$, a clear win at $d=7$ ($+16\%$ over the classic block at
$p=0.003$), and no help at the $d=9/11$ boundary. We position the work
explicitly as a sizing-and-boundary study behind existing LER-beating
decoders---not a scaling law and not a new state-of-the-art. All numbers
are emitted by a single strict reproduction script over seeds 42/43/44.
\end{abstract}

% ====================================================================
\section{Introduction}

\paragraph{Recovery from syndrome evidence and a prior.}
The classical question behind this paper family is whether a sender can
transmit a compact syndrome while a receiver reconstructs using a shared
source model. A syndrome restricts the compatible signals, and the prior
determines which alternatives are plausible. The companion voxel study
examines reconstruction quality with additional transmitted side information.

Quantum error correction has a corresponding inference problem with a
different success criterion. Let $e$ denote a physical error history,
$s$ its detector record, and $\ell(e)$ the logical observable flip.
The noise model gives $P(e)$, and optimal binary logical prediction uses
\[
 \hat\ell(s)=\arg\max_{b\in\{0,1\}}
 \sum_{e:\,s(e)=s,\,\ell(e)=b} P(e).
\]
Thus logical probabilities sum over compatible physical errors. Predicting
a likely individual error, producing a plausible classical reconstruction,
and preserving an unknown quantum state are distinct objectives.
Our neural decoder approximates the binary logical decision from detector
data; it does not construct a physical Hamiltonian or an autonomous memory.

Existing neural and hybrid decoders already improve on particular MWPM
baselines~\cite{bausch2024alphaqubit,lange2025gnn,lee2025mamba,peled2026nmwpm}.
Here we ask a narrower empirical question: how do network architecture,
capacity, and training budget affect a direct 3D convolutional decoder
as code distance increases? We characterize the configurations and
resource budgets actually tested, rather than deriving an intrinsic
capacity requirement for the code.

\paragraph{Contribution.}
We give an empirical, multi-seed characterisation of how a pure 3D
convolutional surface-code decoder performs as a function of
code distance under the tested resource budgets. Concretely:
(i) a per-distance capacity knee in $\textrm{width}\times\textrm{depth}$,
selected on validation by a predeclared rule and reported with three-seed
($42/43/44$) test-set error rates and 95\% confidence intervals; (ii) the
resulting operating map---MWPM-beating at $d=3,5,7$, a high-variance
boundary at $d=9$ under the original budget, improved by the larger-data
follow-up, and a below-baseline regime at $d=11$; and (iii) an
architecture comparison (the classic ResNet block~\cite{he2016resnet}
versus an LC0-style pre-activation, dual-output squeeze-excite,
FiLM-conditioned, ReZero block) showing that block design is
distance-dependent, not second-order. We deliberately do \emph{not} fit a
closed-form scaling law. With $d$ measurement rounds the detector volume grows as $d^{3}$,
increasing the input size and computation at fixed network width, but the multi-seed data
show a budget-dependent \emph{boundary}---a regime where seed variance grows and
the pure-spatial decoder falls below MWPM---rather than a clean power law,
and we report it as such.

\paragraph{Positioning against prior art.}
``Neural decoder beats MWPM'' is not the contribution of this paper.
Lange et al.~\cite{lange2025gnn} beat MWPM on the same rotated surface
code under circuit-level depolarising noise at the same sub-threshold
operating points we use, with larger margins than the pure-CNN decoder
presented here achieves at matched distances. AlphaQubit~\cite{bausch2024alphaqubit}
beats MWPM on real hardware data, and NMWPM~\cite{peled2026nmwpm} beats
MWPM near threshold. Closest to our architecture, Gicev et
al.~\cite{gicev2025fcn} study fully-convolutional 3D networks on the
same spatiotemporal detector volume (including a generative variant) and
generalise a fixed-size, translation-invariant model across distance;
they do not study a capacity-versus-distance sizing rule, which is exactly
the gap we fill. What the present work adds is therefore orthogonal to the
logical-error-rate race: (i) an empirical, multi-seed capacity-sizing map
for a pure spatial decoder as a function of code distance; (ii) a characterization of competitiveness at the tested budgets,
including the $d=9$ improvement with additional training data; and (iii) an architecture comparison testing whether
block design matters once capacity is sized---it does, and in a
distance-dependent way.

This paper is one member of a four-paper program on learned recovery over
constraint systems. The companion occupancy-grid compression
work~\cite{guarraci2026compression} motivates local, prior-aware
refinement from comparisons of specific check-and-solver combinations.
It does not establish a universal locality condition. The gauge-structure note shows the complementary warning that an
exact algebraic characterisation of a codespace need not prescribe a useful
architecture. Surface-code QEC provides the positive-control case for locality:
the stabilizer geometry supplies local, overlapping checks by construction.
The present paper therefore asks a different question: once locality is
guaranteed by physics, how much neural capacity is needed to exploit it as
code distance grows?

\paragraph{What this paper is not.}
We are explicit about the scope so that reviewers can calibrate the
contribution:
\begin{itemize}[nosep]
  \item \textbf{Not autonomous QEC.} A classical network predicts logical
        flips from simulated detector records. No physical relaxation or
        autonomous recovery process is implemented.
  \item \textbf{Not a closed-form scaling law.} We report an empirical
        per-distance capacity knee and a boundary, not a fitted
        $\alpha \cdot d^{3}$ law. The $d^{3}$ detector volume describes input and compute growth;
        it does not imply a lower bound on trainable parameter count; a first-principles capacity bound is a separate problem.
  \item \textbf{Not real hardware.} All syndromes are generated by
        Stim~\cite{gidney2021stim} under circuit-level depolarising
        noise. We do not test on biased noise, leakage, crosstalk, or
        real-device syndromes.
  \item \textbf{Not real-time.} Inference is measured on commodity GPUs
        and remains roughly two orders of magnitude too slow for the
        syndrome cycle time of current superconducting hardware. The
        $O(d^{3})$ scaling and regularity of the 3D convolutional
        kernel make the architecture FPGA/ASIC friendly, but we have
        not built that deployment.
\end{itemize}

\paragraph{Roadmap.}
Section~\ref{sec:background} summarises the rotated
surface code, detector events, and the MWPM baseline. Section~3
describes the 3D convolutional decoder and the error-rate conditioning
that lets each model cover the sampled range of operating points.
Section~4 presents the per-distance capacity knees, the three-seed
test-set results with confidence intervals, the resulting
MWPM-beating / variable / below-baseline operating map, and the
classic-versus-LC0 comparison. Section~5 discusses follow-ups the results
motivate: physics-informed priors that bake code geometry into the decoder
before any data is seen, and the LC0 component ablation and paired
LC0-versus-classic test that the distance-dependent block-design finding
calls for.

% ====================================================================
\section{Background}
\label{sec:background}

\subsection{Surface code basics}

The rotated surface code~\cite{bravyi1998,fowler2012} at distance $d$
encodes one logical qubit in $d^2$ data qubits with $d^2 - 1$
stabilizers arranged in a checkerboard on a 2D grid. Stabilizers are
weight-4 in the bulk and weight-2 at the boundary. $X$ errors flip $Z$
stabilizers and vice versa, so $X$ and $Z$ decoding decouple under
simple noise models. The code has distance $d$, meaning any logical
error requires at least $d$ physical errors, and its error suppression
improves exponentially with $d$ below the code's threshold error
rate~\cite{dennis2002,fowler2012}.

\subsection{Detectors and detector events}

In modern descriptions following Gidney~\cite{gidney2021stim}, syndrome
processing is described in terms of \emph{detectors}---parities of pairs
of consecutive stabilizer measurements that should be equal in the
absence of errors. A flipped detector indicates that an error occurred
between the two measurements. For distance $d$ with $d$ rounds of
measurement, the number of detectors is approximately
$d \times (d^2 - 1)$. For each shot, an ``observable flip'' indicates
whether the encoded logical qubit has been flipped relative to its
initial state.

\subsection{MWPM baseline}

We use \textbf{PyMatching v2}~\cite{higgott2022pymatching,higgott2025sparse},
the standard fast MWPM implementation, configured from Stim's detector
error model with \texttt{decompose\_errors=True}. Comparisons refer to this specific detector-error-model-based
configuration, not every matching variant or tuning. In particular,
decomposing errors into graphlike components does not establish that all
circuit-noise correlations are retained by the matching objective.

% ====================================================================
\section{Method}

\subsection{Spatial embedding of detector events}

Stim~\cite{gidney2021stim} assigns each detector a coordinate tuple
$(x, y, t)$ where $x, y$ are the spatial position of the stabilizer on
the code lattice and $t$ is the measurement round. We extract these
coordinates from the circuit and discretize each axis to integer indices,
yielding a grid size $(D_x, D_y, D_t)$. For the rotated surface code
with $d$ rounds at distance $d$, this gives a $(d{+}1) \times (d{+}1)
\times (d{+}1)$ grid (sparse, checkerboard occupancy).

For each shot, we construct a sparse binary volume
$V \in \{0,1\}^{D_x \times D_y \times D_t}$ where $V[i, j, k] = 1$ if
and only if the detector at coordinate $(i, j, k)$ fired. Cells without
an associated detector are zero.

\subsection{Architecture}

The decoder is a 3D residual convolutional network based on the ResNet
architecture~\cite{he2016resnet} adapted to volumetric data:
\begin{itemize}[nosep]
  \item \textbf{Input convolution} ($1 \to W$) lifts the sparse binary
        volume to a feature volume of constant width $W$
  \item \textbf{$N$ residual blocks}, each containing two \texttt{Conv3d}
        layers, group normalization~\cite{wu2018groupnorm}, SiLU
        activation, and a residual skip connection
  \item \textbf{Error rate conditioning}: a small MLP maps the scalar
        physical error rate $p$ to a conditioning vector that is added to
        the per-block feature representations, analogous to the timestep
        embedding used in diffusion models~\cite{ho2020ddpm}
  \item \textbf{Global average pooling} over spatial dimensions
  \item \textbf{Output MLP head} producing a single logit for logical
        flip prediction
\end{itemize}
The historical fixed-size prototype used $W = 128$ and $N = 12$
(10.7M parameters). The active paper uses per-distance knee cells selected
by the clean capacity scan. The follow-up data-budget experiment separates
one source of performance limitation from parameter count alone.

\subsection{Mixed error rate training}

For each distance and architecture, \textbf{a single model is
trained on a uniform mixture of physical error rates}. For each training
iteration, samples are drawn from circuits with
$p \in \{0.001, 0.002, \ldots, 0.007\}$, and each sample is paired with
its true $p$ as conditioning input. The model learns to adapt its
decoding behavior to the operating point. This is the same approach we
use for variable-rate compression in our companion
paper~\cite{guarraci2026compression}.

\subsection{Coset equivalence}

Two error patterns that differ by a stabilizer element produce the same
syndrome and the same logical effect~\cite{dennis2002,fowler2012}. The
decoder need only identify the coset of the error in the quotient by the
stabilizer group. Because our decoder predicts the logical observable
flip directly (a single binary outcome) rather than the full error
pattern, stabilizer-equivalent errors automatically receive identical
training labels. The decoder is therefore invariant to coset choice by
construction, without requiring explicit augmentation. This is the same
property that other binary-classification QEC decoders rely on, including
AlphaQubit~\cite{bausch2024alphaqubit}.

\subsection{Training details}

We train with AdamW~\cite{loshchilov2019adamw}, learning rate $10^{-3}$,
cosine annealing~\cite{loshchilov2017sgdr} over 30 epochs. Loss is
binary cross-entropy between predicted and true observable flip. Training
data consists of 200K samples per error rate $\times$ 7 rates = 1.4M
total, generated from Stim~\cite{gidney2021stim} circuit-level
depolarizing circuits. Batch size is 2,048. Multi-seed experiments use
seeds 42, 43, 44.

% ====================================================================
\section{Experiments}

\subsection{Setup}

We evaluate on the rotated surface code at distances
$d \in \{3, 5, 7, 9, 11\}$ with circuit-level depolarizing noise.
Training and validation circuits are generated using
Stim~\cite{gidney2021stim} with all four common noise sources: data
depolarization before each round, Clifford depolarization after each
gate, reset bit flip errors, and measurement bit flip errors, all at the
same physical error rate $p$. MWPM baselines are computed using
PyMatching v2~\cite{higgott2022pymatching,higgott2025sparse} with
\texttt{decompose\_errors=True} on the same detector events, enabling
paired significance testing via McNemar's test.

The active clean experiment uses the smallest per-distance knee cell that
is within the predeclared tolerance of the best seed-42 capacity-scan
validation result. Each reported cell must have seeds 42, 43, and 44.
Validation uses \textbf{1M shots per error rate} per seed, sampled
independently from training data.

Data provenance: all results are from k8s Jobs on the galatheus cluster.
The runs used RTX 3090, RTX 4070~Ti, and RTX 4090 GPUs. Job logs are
archived on the shared scratch PVC under \path{/work/qec_runs/}.

\subsection{Per-distance capacity knees and multi-seed results}
\label{sec:results}

For each distance we select the smallest $\textrm{width}\times\textrm{depth}$
cell whose seed-42 validation mean logical error rate (LER) is within a
predeclared tolerance of the best cell in the capacity scan (the
``knee''). The selection rule---minimum mean validation LER across the four
eval rates, equally weighted---is fixed before any test-set evaluation. The
selected knees are listed in Table~\ref{tab:knees}. Every reported cell is
then evaluated on a \emph{disjoint} test split (sampler seed
$\textrm{seed}\times1000+800$, disjoint from train and validation) at all
three seeds $42/43/44$. All numbers below are emitted by
\texttt{reproduce\_qec\_results.py} from the per-run
\texttt{eval\_test\_split.py} JSONs; the reproducer fails on incomplete
seed coverage by construction.

\begin{table}[htbp]
\centering
\caption{Selected per-distance capacity knees
($\textrm{width}\times\textrm{depth}$), each evaluated at seeds $42/43/44$
on the disjoint test split.}
\label{tab:knees}
\begin{tabular}{@{}lccccc@{}}
\toprule
Block & $d{=}3$ & $d{=}5$ & $d{=}7$ & $d{=}9$ & $d{=}11$ \\
\midrule
Classic & $16{\times}4$ & $64{\times}8$ & $64{\times}8$ & $96{\times}12$ & $96{\times}12$ \\
LC0     & $16{\times}4$ & $64{\times}8$ & $96{\times}12$ & $128{\times}12$ & $96{\times}12$ \\
\bottomrule
\end{tabular}
\end{table}

\paragraph{Operating map: win, boundary, failure.}
Table~\ref{tab:boundary} reports the headline at $p=0.003$: the three-seed
mean test LER for each block against the paired MWPM baseline on the same
shots. The pure spatial decoder beats MWPM at $d=3,5,7$, is statistically
tied at $d=9$ (the seed-level $95\%$ CI on the classic LER,
$[0.00045,0.00072]$, straddles the MWPM rate $0.000557$), and falls well
below MWPM at $d=11$ at the selected knee. The $d=9$ and $d=11$ rows are
dominated by seed variance: the per-seed classic $d=9$ LERs are
$\{0.00059,0.00045,0.00072\}$ (one seed beats MWPM, two do not), and the
LC0 $d=9$ LERs are $\{0.00060,0.00049,0.00139\}$ with a $2$--$3\times$
outlier. This is the signature of a budget-dependent performance boundary, not a smooth
power-law trend---which is why we report a boundary rather than fit a law.

\paragraph{The $d=9$ boundary is under-resourcing, not a fundamental limit.}
That boundary is a property of the \emph{selected knee} (the predeclared
width/depth), not of the distance. A follow-up capacity probe at $d=9$---same
architecture (classic, $w{=}96$, depth $12$) but $5\times$ the training
data---\emph{beats} MWPM across the board on the disjoint test split over three
seeds $\{42,43,44\}$: three-seed mean neural LER $0.000232$ vs.\ MWPM $0.000537$
at $p=0.003$ ($+56.5\%$; per-seed McNemar $p\!\sim\!10^{-7}$--$10^{-11}$),
$0.00346$ vs.\ $0.006785$ at $p=0.005$ ($+49\%$), and $0.0201$ vs.\ $0.0312$ at
$p=0.007$ ($+36\%$), with tight cross-seed agreement. So the $d=9$ ``boundary''
of Table~\ref{tab:boundary} is an under-training/under-data artifact of the knee,
not a fundamental capacity wall---a sizing problem more data resolves. (Numbers
emitted by \texttt{reproduce\_qec\_results.py} on the \texttt{qec-bnd-may30}
probe; \texttt{data/d9\_dataarm\_reproduction.json}, see \texttt{REPRODUCIBILITY.md}.)

\begin{table}[htbp]
\centering
\caption{Operating map at $p=0.003$: three-seed mean test-set LER vs.\ the
paired MWPM baseline (same shots), with relative improvement (positive
means the neural decoder beats MWPM). ``Improv.'' is the seed-mean of the
per-seed relative improvement, which at the high-variance $d=9,11$ boundary
differs from the ratio of the mean LERs. Per-seed McNemar tests favour the
neural decoder at $d\le7$ and MWPM at $d=11$ (per-seed $p$-values in the
reproduction JSON).}
\label{tab:boundary}
\begin{tabular}{@{}rrrrrr@{}}
\toprule
 & & \multicolumn{2}{c}{Classic} & \multicolumn{2}{c}{LC0} \\
\cmidrule(lr){3-4}\cmidrule(lr){5-6}
$d$ & MWPM LER & Neural LER & Improv. & Neural LER & Improv. \\
\midrule
3  & 0.006587 & 0.005681 & $+13.7\%$ & 0.005657 & $+14.1\%$ \\
5  & 0.003290 & 0.002734 & $+16.9\%$ & 0.002667 & $+18.9\%$ \\
7  & 0.001396 & 0.001237 & $+11.4\%$ & 0.001037 & $+25.6\%$ \\
9  & 0.000557 & 0.000588 & $-5.6\%$  & 0.000827 & $-49.3\%$ \\
11 & 0.000203 & 0.001049 & $-453\%$  & 0.002455 & $-1192\%$ \\
\bottomrule
\end{tabular}
\end{table}

\paragraph{Block design is distance-dependent, not second-order.}
The two blocks are indistinguishable at $d=3,5$ (LC0 is $+0.4\%$ and
$+2.4\%$ over classic at $p=0.003$), but at $d=7$ the LC0 block is a clear
win: $+16.1\%$ over the classic block, i.e.\ $+25.6\%$ vs.\ MWPM against
the classic block's $+11.4\%$. At the $d=9/11$ boundary LC0 does not help
and is in fact worse than the classic block. The architectural choice is
therefore not second-order: it extends the useful regime at intermediate
distance and offers no rescue at the boundary. This is an aggregate
(unpaired) comparison---the two blocks were trained on different shots; a
paired same-shot McNemar test between blocks
(\texttt{eval\_paired\_comparison.py}) is future work.

\paragraph{Full per-rate results.}
Table~\ref{tab:ler-full} gives the three-seed mean test LER ($\pm$SEM) for
both blocks across all four physical error rates. The qualitative picture
is rate-stable---the win at $d\le7$, the $d=9$ boundary, and the $d=11$
collapse hold across $p\in\{0.001,0.003,0.005,0.007\}$---and the SEM at
$d=9,11$ is an order of magnitude larger than at $d\le7$, quantifying the
boundary's instability.

\begin{table}[htbp]
\centering
\caption{Three-seed mean test-set LER ($\pm$SEM) at the selected knees
(seeds $42/43/44$; disjoint test split).}
\label{tab:ler-full}
\small
\begin{tabular}{@{}llrrrr@{}}
\toprule
Block & $d$ & $p{=}0.001$ & $p{=}0.003$ & $p{=}0.005$ & $p{=}0.007$ \\
\midrule
Classic & 3  & $0.00064{\pm}0.00001$ & $0.00568{\pm}0.00002$ & $0.01543{\pm}0.00005$ & $0.02926{\pm}0.00006$ \\
Classic & 5  & $0.00010{\pm}0.00000$ & $0.00273{\pm}0.00003$ & $0.01276{\pm}0.00021$ & $0.03413{\pm}0.00060$ \\
Classic & 7  & $0.00001{\pm}0.00000$ & $0.00124{\pm}0.00005$ & $0.01002{\pm}0.00041$ & $0.03615{\pm}0.00100$ \\
Classic & 9  & $0.00000{\pm}0.00000$ & $0.00059{\pm}0.00008$ & $0.00769{\pm}0.00041$ & $0.03793{\pm}0.00131$ \\
Classic & 11 & $0.00001{\pm}0.00001$ & $0.00105{\pm}0.00042$ & $0.01235{\pm}0.00335$ & $0.06007{\pm}0.01058$ \\
\midrule
LC0 & 3  & $0.00065{\pm}0.00001$ & $0.00566{\pm}0.00003$ & $0.01541{\pm}0.00004$ & $0.02912{\pm}0.00022$ \\
LC0 & 5  & $0.00010{\pm}0.00000$ & $0.00267{\pm}0.00004$ & $0.01276{\pm}0.00015$ & $0.03398{\pm}0.00027$ \\
LC0 & 7  & $0.00001{\pm}0.00000$ & $0.00104{\pm}0.00002$ & $0.00881{\pm}0.00016$ & $0.03287{\pm}0.00057$ \\
LC0 & 9  & $0.00000{\pm}0.00000$ & $0.00083{\pm}0.00028$ & $0.00964{\pm}0.00242$ & $0.04406{\pm}0.00775$ \\
LC0 & 11 & $0.00003{\pm}0.00001$ & $0.00246{\pm}0.00077$ & $0.02314{\pm}0.00512$ & $0.09092{\pm}0.01301$ \\
\bottomrule
\end{tabular}
\end{table}

\subsection{Inference latency and scaling}

For practical QEC deployment, decoder latency is as important as
accuracy: superconducting hardware produces syndromes at microsecond
rates, and the decoder must keep up with the cycle time to enable
real-time fault-tolerant computation.

A key advantage of our spatially-embedded convolutional approach over
attention-based decoders is its scaling behavior with code distance
(Table~\ref{tab:complexity}).

\begin{table}[htbp]
\centering
\caption{Asymptotic per-shot inference cost vs.\ code distance.}
\label{tab:complexity}
\begin{tabular}{@{}llll@{}}
\toprule
Decoder & Architecture & Per-shot cost & Notes \\
\midrule
MWPM & matching & ${\sim}O(d^3)$ & average case \\
AlphaQubit & recurrent transformer & $O(d^4)$ & attention quadratic in $d^2$ \\
Mamba & state-space & $O(d^3)$ & linear per round \\
\textbf{Ours} & \textbf{3D ResNet} & $\bm{O(d^3)}$ & \textbf{3D conv over $(d{+}1)^3$} \\
\bottomrule
\end{tabular}
\end{table}

Our cost matches MWPM and Mamba and is \textbf{strictly better than
AlphaQubit by a factor of $d$}. At $d=11$ this is already an
order-of-magnitude difference; at $d=21$---the distance currently
targeted by industry roadmaps---the gap is approximately $21\times$.

We measured inference cost on an NVIDIA RTX 4070~Ti
(Table~\ref{tab:latency}).

\begin{table}[htbp]
\centering
\caption{Measured inference latency (width 128, depth 12).}
\label{tab:latency}
\begin{tabular}{@{}rrrrrr@{}}
\toprule
$d$ & Grid & Detectors & Batched ($\mu$s/shot) & Shots/sec & Single ($\mu$s) \\
\midrule
3 & $4^3$ & 24 & 49 & 20{,}243 & 2{,}371 \\
5 & $6^3$ & 120 & 181 & 5{,}528 & 2{,}462 \\
7 & $8^3$ & 336 & 429 & 2{,}333 & 2{,}324 \\
9 & $10^3$ & 720 & 835 & 1{,}198 & 2{,}344 \\
11 & $12^3$ & 1{,}320 & 1{,}441 & 694 & 3{,}001 \\
\bottomrule
\end{tabular}
\end{table}

The batched latency scales as ${\sim}d^{2.9}$, slightly better than the
theoretical $d^3$ because small distances do not saturate GPU compute.
Single-shot latency is approximately constant at 2.3--3.0\,ms, dominated
by GPU kernel launch overhead.

\textbf{Deployment scope.} These timings describe the fixed width-128,
depth-12 prototype in Table~\ref{tab:latency}, not every selected
configuration. Batched throughput and single-shot response latency are
different deployment constraints. Accelerator compilation, distillation,
quantization, and dedicated hardware are possible directions, but no
combined speedup or real-time deployment is demonstrated here.

\subsection{Comparison with prior work}

We do not claim state-of-the-art logical error rates, and the comparison to
prior neural decoders is deliberately at the level of regime rather than a
single headline LER. AlphaQubit~\cite{bausch2024alphaqubit}, Lange et
al.~\cite{lange2025gnn}, Mamba-style state-space
decoders~\cite{lee2025mamba}, and NMWPM~\cite{peled2026nmwpm} already beat
MWPM in their respective settings, and Lange et al.\ do so by larger
relative margins than the pure-CNN decoder here achieves at matched
distance. Our contribution is orthogonal: the per-distance sizing map, the
boundary location, and the distance-dependent block-design effect. The
fully-convolutional 3D decoders of Gicev et al.~\cite{gicev2025fcn}
are the closest architecture; they rely on a fixed-size,
translation-invariant network to generalise across distance, whereas we
quantify the capacity a non-shared per-distance model needs and the
distance at which even that runs out.

\subsection{Limitations}

We are explicit about what we have and have not shown:
\begin{itemize}[nosep]
  \item All experiments are on \textbf{simulated} circuit-level
        depolarizing noise. We have not tested on real hardware data,
        biased noise, leakage, or crosstalk.
  \item The $d=9$ and $d=11$ results have large seed-to-seed variance
        (SEM an order of magnitude larger than at $d\le7$); the boundary
        is reported as such, and the $d=11$ knee is explicitly
        a below-baseline result at the tested budget, rather than a tuned best case.
  \item The classic-versus-LC0 comparison is aggregate (unpaired): no
        paired per-shot comparison between the two trained models is
        reported. Different training shots do not preclude evaluating
        both on a common test set; that comparison remains a follow-up.
  \item The LC0 win at $d=7$ is not decomposed into its components
        (pre-activation, dual-output squeeze-excite, FiLM, ReZero); a
        component ablation is future work.
  \item The larger capacity-scan directories used to select the retained
        configurations are not included in the public release. The
        released evaluation JSONs reproduce their reported performance,
        but do not independently reproduce the full selection process.
  \item AlphaQubit's headline 25\% result is on real hardware (Sycamore
        data not publicly available); our simulation result is not
        directly comparable.
\end{itemize}

% ====================================================================
\section{Connection to Constraint Locality}

Our companion compression paper investigates constraint locality in
a setting where the parity checks are a design choice rather than a physical
given~\cite{guarraci2026compression}. In that domain, fixed high-weight
projection can damage a strong learned proposal, while local overlapping
checks can provide useful refinement signals.

Surface-code decoding gives the corresponding positive-control case.
Stabilizers have exactly the required structure \emph{by construction}: each
stabilizer acts on at most four adjacent qubits (weight-4 in the bulk,
weight-2 at the boundary), neighboring stabilizers share support on data
qubits, and a single Pauli error at a fixed time has a local syndrome footprint.
Circuit faults can generate correlated detector events, so this geometric
observation does not reduce every circuit-level event to an independent
pair. The
compression problem must engineer locality; the surface code supplies it
through its geometry.

This connection motivates the present capacity-scaling question. If local
overlapping constraints are the regime where learned syndrome decoders can
use spatial structure, then a surface-code decoder should expose the detector
geometry to the network rather than flattening the syndrome. Our 3D
convolutional embedding does exactly that. The remaining question is not
whether locality is present, but how much neural capacity is needed to
exploit it as the detector volume grows with distance.

% ====================================================================
\section{Discussion and Future Work}

\paragraph{The boundary depends on the training budget.}
The original selected configurations beat the implemented MWPM baseline
through $d=7$, are variable at $d=9$, and underperform at $d=11$.
The $d=9$ follow-up uses the same classic architecture with five times
the training data and achieves lower error than MWPM across its reported
rates. The original boundary therefore is not a demonstrated limit of
that architecture. At $d=11$, the existing experiments do not isolate
parameter capacity, data budget, and optimization as causes of the gap.
We report a budget-dependent operating map, not a scaling law or a
lower bound on the capacity needed to decode.

\paragraph{Block design is distance-dependent.}
We had hypothesised that, once capacity is sized, the residual block design
would be second-order---that the capacity requirement is a property of the
decoding problem rather than of the network. The multi-seed data refute the
strong version of this. The LC0 block is neutral at $d=3,5$, a clear win at
$d=7$ ($+16\%$ over the classic block), and no help---indeed worse---at the
$d=9/11$ boundary. Block design matters, and matters most at intermediate
distance. Which LC0 component carries the $d=7$ win (pre-activation,
dual-output squeeze-excite, FiLM conditioning, or ReZero) is the most
immediate follow-up; it requires block variants that isolate each component
and a paired same-shot test against the classic block.

\paragraph{$d=11$ remains below the baseline at the tested budget.}
At the $d=11$ selected configuration ($w=96$, depth $12$), the
three-seed classic LER at $p=0.003$ is $0.00105$, versus MWPM's
$0.00020$, with LC0 performing worse. These observations warrant
controlled capacity and data-budget sweeps; they do not establish an
intrinsic stopping distance for a pure convolutional decoder.

\paragraph{The most promising follow-up: a physics-informed decoder.}
One hypothesis is that explicit code structure could reduce the data or
capacity required by a pure-spatial decoder. The present experiments do
not measure how much capacity is spent learning that structure. A
physics-informed decoder parameterises $H(x) = H_{\text{code}}(x) +
H_{\theta}(x)$, with $H_{\text{code}}$ fixed from the code (a
differentiable form of the MWPM edge weights) and only the residual
$H_{\theta}$ learned. This is the direction in which hybrid decoders
already beat MWPM more decisively than pure-neural ones: Lange et
al.~\cite{lange2025gnn} and especially NMWPM~\cite{peled2026nmwpm}, which
learns MWPM edge weights and reports near distance-independent parameter
efficiency. Because NMWPM already realises much of this idea, the novelty
of a physics-informed variant must be argued against it directly rather
than against MWPM. The operating map identifies configurations where a hybrid comparison
would be informative. Such a comparison should include the larger-data
classic decoder, since additional data already improves the $d=9$ result.

\paragraph{Generative and score-matching decoders.}
A score- or diffusion-style head that emits a per-cell gradient rather than
a single logit is a natural denser-signal alternative, with a test-time
compute knob via iterative refinement~\cite{ho2020ddpm}. This direction is
already active: diffusion and masked-diffusion QEC decoders, and the
generative variant of the fully-convolutional decoder of Gicev et
al.~\cite{gicev2025fcn}, appeared concurrently with this work. We
therefore flag it as relevant prior art rather than as our own
contribution.

\paragraph{Realistic noise and deployment.}
All experiments use simulated circuit-level depolarizing noise.
Performance under biased noise, leakage, crosstalk, and real-device
distribution shift remains unmeasured. Inference latency also remains
above the intended real-time regime. Quantization, accelerator compilation,
and dedicated hardware are possible follow-ups, but their speedups and
effects on logical error have not been demonstrated here. These are
additional experimental questions, separate from the reported resource map.

\section{Conclusion}
We presented a three-seed study of direct 3D convolutional surface-code
decoders under specified architecture and training budgets. In the
original selected configurations, the decoder improves on the implemented
MWPM baseline at $d=3,5,7$, is variable at $d=9$, and underperforms at
$d=11$. A fivefold-data follow-up at $d=9$ improves on MWPM without
changing the classic architecture, demonstrating that the original
boundary is budget-dependent. The LC0-style block helps at $d=7$ in
the retained comparison; its component contributions and paired
performance against the classic block remain to be tested.

The contribution is an empirical operating map, not a scaling law or
a state-of-the-art decoding claim. It addresses the paper family's
central inference question in a quantum setting: given syndrome evidence
and a prior over errors, which learned representation and resource
budget suffice to infer the relevant logical information? Source code
and retained evaluation JSONs are released; re-evaluation of trained
models requires the external checkpoints described in the reproduction
guide.

% ====================================================================

\bibliographystyle{unsrt}

\begin{thebibliography}{25}

\bibitem{kitaev2003}
A.~Yu.~Kitaev, ``Fault-tolerant quantum computation by anyons,''
\emph{Ann.\ Phys.}, vol.~303, no.~1, pp.~2--30, 2003.

\bibitem{bravyi1998}
S.~B.~Bravyi and A.~Yu.~Kitaev, ``Quantum codes on a lattice with
boundary,'' arXiv:quant-ph/9811052, 1998.

\bibitem{dennis2002}
E.~Dennis, A.~Kitaev, A.~Landahl, and J.~Preskill, ``Topological
quantum memory,'' \emph{J.~Math.\ Phys.}, vol.~43, no.~9,
pp.~4452--4505, 2002.

\bibitem{fowler2012}
A.~G.~Fowler, M.~Mariantoni, J.~M.~Martinis, and A.~N.~Cleland,
``Surface codes: Towards practical large-scale quantum computation,''
\emph{Phys.\ Rev.\ A}, vol.~86, p.~032324, 2012.

\bibitem{google2023suppressing}
R.~Acharya et~al.\ (Google Quantum AI), ``Suppressing quantum errors by
scaling a surface code logical qubit,'' \emph{Nature}, vol.~614,
pp.~676--681, 2023.

\bibitem{google2024below}
R.~Acharya et~al.\ (Google Quantum AI), ``Quantum error correction below
the surface code threshold,'' \emph{Nature}, vol.~638, pp.~920--926,
2025.

\bibitem{higgott2022pymatching}
O.~Higgott, ``PyMatching: A Python package for decoding quantum codes
with minimum-weight perfect matching,'' \emph{ACM Trans.\ Quantum
Computing}, vol.~3, no.~3, pp.~1--16, 2022.

\bibitem{higgott2025sparse}
O.~Higgott and C.~Gidney, ``Sparse Blossom: correcting a million errors
per core second with minimum-weight matching,'' \emph{Quantum}, vol.~9,
p.~1600, 2025.

\bibitem{peled2026nmwpm}
N.~Peled, S.~Zenati, and E.~Nachmani, ``Neural minimum-weight perfect
matching,'' arXiv:2601.00242, 2026.

\bibitem{delfosse2021}
N.~Delfosse and N.~H.~Nickerson, ``Almost-linear time decoding
algorithm for topological codes,'' \emph{Quantum}, vol.~5, p.~595, 2021.

\bibitem{bausch2024alphaqubit}
J.~Bausch et~al., ``Learning high-accuracy error decoding for quantum
processors,'' \emph{Nature}, 2024. (AlphaQubit.)

\bibitem{lee2025mamba}
C.~Lee, T.~Hur, and D.~K.~Park, ``Scalable neural decoders for
practical real-time quantum error correction,'' arXiv:2510.22724, 2025.

\bibitem{torlai2017}
G.~Torlai and R.~G.~Melko, ``Neural decoder for topological codes,''
\emph{Phys.\ Rev.\ Lett.}, vol.~119, p.~030501, 2017.

\bibitem{varsamopoulos2018}
S.~Varsamopoulos, B.~Criger, and K.~Bertels, ``Decoding small surface
codes with feedforward neural networks,'' \emph{Quantum Sci.\ Technol.},
vol.~3, no.~1, p.~015004, 2018.

\bibitem{baireuther2018}
P.~Baireuther et~al., ``Machine-learning-assisted correction of
correlated qubit errors in a topological code,'' \emph{Quantum}, vol.~2,
p.~48, 2018.

\bibitem{lange2025gnn}
M.~Lange et~al., ``Data-driven decoding of quantum error correcting
codes using graph neural networks,'' \emph{Phys.\ Rev.\ Research},
vol.~7, p.~023181, 2025.

\bibitem{gicev2025fcn}
S.~Gicev, L.~C.~L.~Hollenberg, and M.~Usman, ``Fully convolutional 3D
neural network decoders for surface codes with syndrome circuit noise,''
arXiv:2506.16113, 2025.

\bibitem{gidney2021stim}
C.~Gidney, ``Stim: a fast stabilizer circuit simulator,''
\emph{Quantum}, vol.~5, p.~497, 2021.

\bibitem{he2016resnet}
K.~He, X.~Zhang, S.~Ren, and J.~Sun, ``Deep residual learning for
image recognition,'' in \emph{CVPR}, pp.~770--778, 2016.

\bibitem{wu2018groupnorm}
Y.~Wu and K.~He, ``Group normalization,'' in \emph{ECCV}, pp.~3--19,
2018.

\bibitem{ho2020ddpm}
J.~Ho, A.~Jain, and P.~Abbeel, ``Denoising diffusion probabilistic
models,'' in \emph{NeurIPS}, vol.~33, pp.~6840--6851, 2020.

\bibitem{loshchilov2019adamw}
I.~Loshchilov and F.~Hutter, ``Decoupled weight decay regularization,''
in \emph{ICLR}, 2019.

\bibitem{loshchilov2017sgdr}
I.~Loshchilov and F.~Hutter, ``SGDR: Stochastic gradient descent with
warm restarts,'' in \emph{ICLR}, 2017.

\bibitem{hu2018squeeze}
J.~Hu, L.~Shen, and G.~Sun, ``Squeeze-and-excitation networks,'' in
\emph{CVPR}, pp.~7132--7141, 2018.

\bibitem{gallager1962ldpc}
R.~G.~Gallager, ``Low-density parity-check codes,'' \emph{IRE Trans.\
Inf.\ Theory}, vol.~8, no.~1, pp.~21--28, 1962.

\bibitem{nachmani2016}
E.~Nachmani, Y.~Be'ery, and D.~Burshtein, ``Learning to decode linear
codes using deep learning,'' in \emph{Allerton Conf.}, pp.~341--346,
2016.

\bibitem{guarraci2026compression}
B.~Guarraci, ``Constraint locality in learned syndrome decoding for 3D
occupancy grids,''
companion paper, 2026.

\end{thebibliography}

\end{document}

